% Source-derived chapter 05: Selmer sheaves and $p$-adic heights
% Original source: universal-p-adic-gross-zagier-formula
\chapter{Selmer sheaves and $p$-adic heights}
\label{universal-p-adic-gross-zagier-formula-chapter-05}
\source{universal-p-adic-gross-zagier-formula}

\begin{remark}[Source section]
\label{universal-p-adic-gross-zagier-formula-chapter-05-source}
\source{universal-p-adic-gross-zagier-formula}
\source{universal-p-adic-gross-zagier-formula}
This chapter reproduces the corresponding section of the retrieved
LaTeX source, including its definitions, statements, examples, and
proofs. It has not yet been translated into Lean.
\notready
\end{remark}

% The source section title was: Selmer sheaves and $p$-adic heights
\label{big sec: sel}
\source{universal-p-adic-gross-zagier-formula}
In this section we present the   theory of Selmer complexes and $p$-adic heights  needed in the rest of the paper. The foundational material  is  taken from the   book of \nek\  \cite{nek-selmer}.
\subsection{Continuous cohomology} 
Let
  $(R^{\circ}, \mathfrak{m})$ be a complete Noetherian local ring,  let $G$ be a topological group. 
  \subsubsection{Continuous cochains for (ind-) admissible $R[G]$-modules} \lb{ssec adm}
Let   $M$  be an $R^{\circ}[G]$-module.
  We say that $M$ is  \emph{admissible of finite type }   if it is of finite type as an $R^{\circ}$-module  and  
   the action $G\times M\to M$ is continuous  (when $M$ is given the $\mathfrak{m}$-adic topology). We say that $M$ is \emph{ind-admissible} if $M=\bigcup_{\alpha}M_{\alpha}$ where $\{M_{\alpha}\}$ is the set of   finite-type admissible $R^{\circ}[G]$-submodules of $M$. 

The complex of continuous cochains of $M$ is denoted by $C^{\bullet}_{\rm cont}(G, M)$; it is defined in the usual  way \cite[(3.4.1)]{nek-selmer} when $M$ is admissible of finite type,  and by  $C^{i}_{\rm cont}(G, M):= \varinjlim_{\alpha} C^{i}_{\rm cont}(G, M_{\alpha})$ when we have a presentation $M=\bigcup_{\alpha}M_{\alpha}$ as above.   The image of  $C^{\bullet}_{\rm cont}(G, M)$ in the derived category of   $D({}_{R}{\rm Mod})$ of $R^{\circ}$-modules is denoted by 
$$ {\rm R}\Gamma(G, M)$$
and its cohomology groups by $$H^{i}(G, M)$$
(we omit the subscript `cont' as we will only be working with continuous cohomology). 
 
\subsubsection{Localisation} Let  
$$R=R^{\circ}[\mathcal{S}^{-1}]$$
 for some multiplicative subset $\mathcal{S}\subset R^{\circ}$, and let $M$ be an $R[G]$-module. We say that $M$ is ind-admissible if it is ind-admissible as an $R^{\circ}[G]$-module, and that it is of finite type if it is of finite type as an $R$-module. Suppose that   $M:=M^{\circ}\otimes_{R^{\circ}} R$ for an ind-admissible $R^{\circ}[G]$-module $M^{\circ}$. Then $M$ is ind-admissible  as an $R^{\circ}[G]$-module and   there is a canonical isomorphism
\beq\label{localisation}
\source{universal-p-adic-gross-zagier-formula}
C^{\bullet}_{\rm cont}(G, M)\cong C^{\bullet}_{\rm cont}(G, M^{\circ})\otimes_{R^{\circ}}R
\eeq
(\cite[(3.7.4)]{nek-selmer}). 


\begin{rema}
\source{universal-p-adic-gross-zagier-formula}
\notready \lb{calCR}Let 
$$\mathscr{C}= \mathscr{C}_{R^{\circ}}$$ 
be the category of schemes isomorphic to open subschemes of $\Spec R^{\circ}$, with maps being open immersions. It follows  from  the previous paragraph that,  for any object $X$ of $\mathscr{C}$,  the condition of ind-admissibility is defined for all quasicoherent $\OO_{X}[G]$-modules, and  the functors ${\rm R}\Gamma(G,- )$  are well-defined on ind-admissible  $\OO_{X}[G]$-modules. Moreover, both the ind-admissibility condition and the functors ${\rm R}\Gamma(G, - )$  are compatible with restriction along open immersions in $\mathscr{C}$. 
\end{rema}

In  the  following, we will not further comment on the generalisation indicated in the previous remark when referring to sources only considering $R^{\circ}[G]$-modules.

\subsubsection{Completed product} \lb{compl product} For $i=1,2$, let $R_{i}^{\circ}$ be complete noetherian local rings, and let $R^{\circ}:= R_{1}^{\circ}\hat{\ot} R_{2}^{\circ}$. We have  a functor 
\beq\lb{hatprod}\hat{\ts}\colon \mathscr{C}_{R_{1}^{\circ}}\ts \mathscr{C}_{R_{2}^{\circ}} \to  \mathscr{C}_{R^{\circ} }\eeq
defined on objects 
by $\Spec R_{1}^{\circ}[1/f_{1}] \hat{\ts} \Spec R_{1}^{\circ}[1/f_{2}]:= \Spec R_{1}^{\circ}\hat{\ot}R_{2}^{\circ}[1/f_{1}\ot1, 1/1\ot f_{2}]$ and glueing.

\subsubsection{Notation} 
Throughout the rest of this section, 
$X$ will denote an object of $\mathscr{C}_{R^{\circ}}$.  If $\calA=\OO_{X}, \OO_{X}[G]$, we denote  by $D({}_{\calA}{\rm Mod})$ the derived category of $\calA$-moduels. We use sub- or superscripts $$\text{ft,  {ind-adm}, $+$, $-$, {b}, $[a, b]$,  {perf},} $$ to denote the full subcategory of objects quasi-isomorphic to complexes of $\calA$-modules that are respectively termwise of finite type, termwise ind-admissible, bounded below, bounded above, bounded, concentrated in degrees $[a,b]$, bounded, perfect (= bounded and termwise projective and of finite type).


\begin{prop}[{ \cite[(3.5.6)]{nek-selmer}}]
\source{universal-p-adic-gross-zagier-formula}
\notready
The functor ${\rm R}\Gamma(G, -)$ can be extended  to a functor on the category of bounded-below complexes of ind-admissible $\OO_{X}[G]$-modules, with values in bounded-below complexes of $\OO_{X}$-modules \cite[(3.4.1.3), (3.5.1.1)]{nek-selmer}. It descends to an exact   functor  
$${\rm R}\Gamma(G, -)\colon D^{+}({}^{\textup{ind-adm}}_{\OO_{X}[G]}{\rm Mod})\to D^{+}({}_{\OO_{X}}{\rm Mod}).$$
\end{prop}
 

\subsubsection{Base-change} Suppose that $R\twoheadrightarrow R'=R/I$ is a surjective map of rings. Let $j\colon R^{\circ} \to R= R^{\circ}[\mathscr{S}^{-1}]$ be the natural map and let  $I^{\circ}:= j^{-1}(I)$. Then $R^{\circ}{}':= R^{\circ}/I^{\circ}$ is also complete local Noetherian, and  we may write $R'= R^{\circ}{}'[\mathscr{S}']^{-1}$ where  $\mathscr{S}'$ is the image of $\mathscr{S}$ in $R^{\circ}{}'$. Let $M'$ be an ind-admissible $R'$-module, then $C^{\bullet}_{\rm cont}(G, M')$ is the same whether we consider $M'$ as an $R'$-module or as an $R$-module: in the special case $R=R^{\circ}$ this follows from the fact that the maximal ideal of $R^{\circ}{}'$ is the image of the maximal ideal of $R^{\circ}$, so that the $\mathfrak{m}$-adic and $\mathfrak{m}'$-adic topologies on finitely generated $R^{\circ}{}'$-modules coincide; the general case   follows from  the special case by localisation \eqref{localisation}. 

More generally, if $Y\subset X $ is a closed subset, the functor ${\rm R}\Gamma(G, -)$ on $\OO_{Y}[G]$-modules coincides with the restriction of the functor on $\OO_{X}[G]$-modules of the same name.

 \begin{prop}
\source{universal-p-adic-gross-zagier-formula}
\notready\label{base change} Let $M$ be an ind-admissible $\OO_{X}[G]$-module and let $N$ be an $\OO_{X}$-module of finite projective dimension. Then  there is a natural isomorphism  in $D^{\rm b}({}_{\OO_{X}}{\rm Mod})$
\source{universal-p-adic-gross-zagier-formula}
 $$  {\rm R}\Gamma(G, M) \stackrel{\rm L}\otimes_{\OO_{X}} N\cong  {\rm R}\Gamma(G, M \stackrel{\rm L}\otimes_{\OO_{X}} N). $$
  \end{prop}
 \begin{proof} Let $P^{\bullet}$ be a finite projective resolution of $N$.   The natural map  of complexes of $\OO_{X}$-modules
$$C_{\rm cont}^{\bullet}(G, M)\otimes_{\OO_{X} }P^{\bullet}\to  C_{\rm cont}^{\bullet}(G, M\otimes_{\OO_{X}}P^{\bullet})$$
 is an isomorphism by   \cite[(3.4.4)]{nek-selmer}.\footnote{In \emph{loc. cit.}, the ring denoted by $R$ is our $R^{\circ}$, but as our $X$ is open in $\Spec R^{\circ}$, the $\OO_{X}$-modules $P^{n}$ are also flat as $R^{\circ}$-modules and the  cited result applies.} The desired result follows from the definition of derived tensor product. 
 \end{proof}
The proposition applies when $N=\OO_{Y}$ with $Y\subset X$  a local complete intersection, 
or when  $X$ is regular and $N$ is any coherent $\OO_{X}$-module.
 We highlight the following case. 
 \begin{coro}   \label{base change cor} 
\source{universal-p-adic-gross-zagier-formula}
 Let $M$ be an ind-admissible $\OO_{X}[G]$-module that is flat as an $\OO_{X}$-module, and let $x\in X$ be a nonsingular point. Then there is an isomorphism in $D^{\rm b}({}_{\kappa(x)}{\rm Mod})$
 $$ {\rm R}\Gamma(G, M) \stackrel{\rm L}\otimes_{\OO_{X}} \kappa(x)\cong 
 {\rm R}\Gamma(G, M \otimes_{\OO_{X}} \kappa(x)) $$
 hence a second-quadrant spectral sequence 
 $${\rm Tor}_{-p}(H^{q}(G, M), \kappa(x)) \Rightarrow  H^{q-p}(G , M\otimes_{\OO_{X}}\kappa(x)).$$
 \end{coro} 
\begin{proof} After possibly localising at $x$, we may assume that $X=\Spec R$ is the spectrum of a  local ring, which by assumption will be regular. Then  
 $\kappa(x)$   has finite projective dimension over $R$, and the result follows from the previous proposition. 
\end{proof}
 
\subsubsection{Continuous cohomology as derived functor} For $i=0, 1$, the functors $M\mapsto H^{i}(G,M)$ on the category of ind-admissible $R$-modules  coincide with the $i^{\rm th}$ derived functors of $M\mapsto M^{G}$ (\cite[(3.6.2)(v)]{nek-selmer}).


\subsection{Specialisations}\label{(F)} 
\source{universal-p-adic-gross-zagier-formula}
From here on we further  assume that $R^{\circ}$ has finite residue field of characteristic~$p$. 

\subsubsection{Finiteness conditions}
Let $G$ be a profinite group. We consider the condition 
$$\text{(F)}\qquad H^{i}(G, M) \text{ is finite for all finite discrete ${\bf F}_{p}[G]$-modules and all $i\geq 0$}$$
and define the $p$-cohomological dimension of $G$ to be
$${\rm cd}_{p}(G):= \sup\, \{i \, :\, \exists \text{ a finite discrete ${\bf F}_{p}[G]$-module $M$ with } H^{i}(G, M)\neq 0\}.$$
 \begin{lemm}\label{finite coh}
\source{universal-p-adic-gross-zagier-formula}
 If $G$ satisfies (F) then the cohomology groups of ind-admissible $\OO_{X}[G]$-modules of finite type are $\OO_{X}$-modules of finite type (\cite[(4.2.5), (4.2.10)]{nek-selmer}).
  The cohomology of any ind-admissible $\OO_{X}[G]$-module vanishes in degrees $>{\rm cd}_{p}(G)$ (\cite[(4.26)]{nek-selmer}). 
 \end{lemm}
 
 When $E$ is a number field, $S$ is a finite set of places of  $E$ and $G=G_{E, S}$,
  condition (F) is satisfied and ${\rm cd}_{p}(G)= 3$. When $E_{w}$ is a local field and $G=G_{E_{w}}$,
   condition (F) is satisfied and ${\rm cd}_{p}(G)= 2 $. In the latter case we use the notation $H^{i}(E_{w}, M) $ for $H^{i}(G, M)$. 

\subsubsection{Projective limits, specialisations} We give two  results on the compatibility of  $G$-cohomology  with other functors.


\begin{lemm}\label{projlim ok} Let $ G$ be a profinite group satisfying  (F) and let $M=\varprojlim_{n} M_{n}$ be the limit of  a countable projective system of admissible $R^{\circ}$-modules of finite type. Then for all $i$ the natural map
\source{universal-p-adic-gross-zagier-formula}
$$ H^{i}(G, M) \to  \varprojlim_{n} H^{i}(G, M_{n})$$
is an isomorphism.
\end{lemm}
\begin{proof}
In the special case $M_{n}=M/\mathfrak{m}^{n}M$, it is shown in \cite[Corollary 4.1.3]{nek-selmer} that the map under consideration is surjective with kernel $\lim^{(1)}_{n}  H^{i-1}(G, M_{n})$; this vanishes since by (F) those cohomology groups are finite, hence the projective system they form satisfies the Mittag-Leffler condition.  The general case follows from applying the special case to $M$ and the $M_{n}=\varprojlim_{r}M_{n}/\mathfrak{m}^{r}M_{n}$. 
\end{proof}

\begin{prop}
\source{universal-p-adic-gross-zagier-formula}
\notready\label{torss} Let $G$ be a profinite group satisfying (F) and ${\rm cd}_{p}(G)=e<\infty$. Let $M$  be an ind-admissible $\OO_{X}[G]$-module of finite type. Let  $x\in X$ be a nonsingular point, let $i_{0}\geq 0$  and suppose that   
\source{universal-p-adic-gross-zagier-formula}
$$H^{i}(G, M\otimes_{R}\kappa(x))=0$$
for all $i\geq i_{0}+1$. 
\begin{enumerate}
\item For all $i\geq i_{0}+1$, the support of the finitely generated $R$-module $H^{i}(G, M)$ is a proper closed subset not containing $x$.
\item The natural map
$$H^{i_{0}}(G, M)\otimes_{\OO_{X}}\kappa(x)\to H^{i_{0}}(G, M\otimes_{R} \kappa(x))$$
is an isomorphism.
\item\lb{torss3} Suppose further  that $i_{0}=1$,  and that for $y$ in some dense open subset of $X$,  $\dim _{\kappa(y)} H^{1}(G, M\ot\kappa(y))= \dim _{\kappa(x)} H^{1}(G, M\ot\kappa(x))$. Then the natural map
$$H^{{0}}(G, M)\otimes_{R}\kappa(x)\to H^{{0}}(G, M\otimes_{\OO_{X}} \kappa(x))$$
is an isomorphism.
\end{enumerate}
\end{prop}
\begin{proof}  By Nakayama's lemma  and the vanishing assumption, the first statement is equivalent to  
\beq \label{dhi} 
\source{universal-p-adic-gross-zagier-formula}
H^{i}(G, M)\otimes_{\OO_{X}}\kappa(x)\cong H^{i}(G, M\otimes_{\OO_{X}} \kappa(x)).
\eeq
Therefore, for the proof of the first and second statements  it is enough to prove  \eqref{dhi}  for all $i\geq i_{0}$, which we do by decreasing induction on $i$.

 For $i\geq e+1$ the result is automatic. In general, Corollary  \ref{base change cor} gives a second-quadrant spectral sequence 
\beq \lb{E2pq}
E_{2}^{p,q}={\rm Tor}_{-p}^{\OO_{X}}(H^{q}(G, M),  \kappa(x))\Rightarrow H^{q-p}(G, M\otimes_{\OO_{X}} \kappa(x)).\eeq
By induction hypothesis, all terms on the diagonal  $q-p=i$ vanish except possibly  the one with $p=0$, and the differentials with source and target such term are $0$. It follows that $H^{i}(G, M\otimes_{\OO_{X} }\kappa(x))=E_{\infty }^{0, i}=E_{2}^{0, i}=H^{i}(G, M)\otimes_{\OO_{X}}\kappa(x) $.

Finally, under the assumptions of part 3, the finitely generated $R$-module $H^{1}(G, M)$ is locally free of constant rank in a neighbourhood of $x$. Hence in the exact sequence 
$$0\to H^{{0}}(G, M)\otimes_{R}\kappa(x)\to H^{0}(G, M\otimes_{\OO_{X}} \kappa(x)) \to {\rm Tor}_{1}^{\OO_{X}}(H^{1}(G, M), \kappa(x))$$
deduced from \eqref{E2pq}, the last term vanishes.
\end{proof}


  
\subsection{Selmer complexes and height pairings}\label{sec: sel} As in the preceding subsection, let $R^{\circ}$ be a Noetherian local ring  with finite residue field  of characteristic $p$,  $X$ an object of $\mathscr{C}_{R}$. 
\source{universal-p-adic-gross-zagier-formula}

When $E_{w}$ is a local field, we write ${\rm R}\Gamma(E_{w}, -):= {\rm R}\Gamma(G_{E_{w}}, -)$ and similarly for its cohomology groups. For number fields, we will only use the analogous shortened notation for Selmer groups.
\subsubsection{Greenberg data} Let $E$ be a number field, $Sp$  a finite set of finite places of $E$ containing those above $p$. 
 Fix for every $w\vert p$ an embedding $\baar{E}\into\baar{ E}_{{w}}$ inducing an embedding $G_{w}:=G_{E_{w}}\into G_{E, Sp}$.
 If $M$ is an $\OO_{X}[G_{E, Sp}]$-module, we   denote by $M_{w}$ the module $M$ considered as an $\OO_{X}[G_{w}]$-module.
 
 \begin{defi}
\source{universal-p-adic-gross-zagier-formula}
\notready\label{grbg datum} A  \emph{Greenberg datum}   $ ( M, (M_{w}^{+})_{w\in  Sp})$ (often abusively abbreviated by $M$ in what follows)  over $X$ consists of
\source{universal-p-adic-gross-zagier-formula}
 \begin{itemize}
 \item an ind-admissible  $\OO_{X}[G_{E, Sp}]$-module $M$,   finite and locally free  as an $\OO_{X}$-module;
 \item  for every $w\in Sp$ a \emph{Greenberg local condition}, that is  a short exact sequence $$ 0 \to M_{w}^{+}\stackrel{i_{w}^{+}}{\to} M_{w}\to M_{w}^{-}\to 0$$
of ind-admissible $\OO_{X}[G_{w}]$-modules,  finite and locally free as $\OO_{X}$-modules. 
 \end{itemize}
 \end{defi}
 In this paper, at places $w \nmid p$ we will only consider the \emph{strict} Greenberg conditions $M_{w}^{+}=0$.  
 















\subsubsection{Selmer complexes} Given a Greenberg  datum  $M= (M, (M_{w}^{+})_{w\in Sp})$,
  the \emph{Selmer complex}  
$$\wtil{{\rm R}{\Gamma}}_{f}(E, M)$$ is the image of the complex
$$ {\rm Cone}\left(     C^{\bullet}_{\rm cont}(G_{E, Sp},M)\oplus\bigoplus_{w\in Sp} C^{\bullet}_{\rm cont}(E_{w}, M_{w}^{+})\stackrel{ \oplus_{w}{\rm res}_{w}-i_{w, *}^{+} }{\longrightarrow}    \bigoplus_{w\in Sp} C^{\bullet}_{\rm cont}(E_{w}, M_{w}) \right)[-1]  $$
  in $D({}_{\OO_{X}}^{\rm ft}{\rm Mod})$. Its cohomology groups are denoted by $\wtil{H}^{i}_{f}(E, M)$. 
We   have an exact triangle 
\beq\label{ex tri}
\source{universal-p-adic-gross-zagier-formula}
\wtil{{\rm R}{\Gamma}}_{f}(E, M)\to {\rm R}{\Gamma}_{}(E, M)\to \oplus_{w\in Sp} {\rm R}\Gamma(E_{w}, M^{-}_{w}). 
 \eeq
  

  
  
  
  
  


\begin{prop}
\source{universal-p-adic-gross-zagier-formula}
\notready The Selmer complex $\wtil{{\rm R}{\Gamma}}_{f}(E, M)$  and all terms of  \eqref{ex tri} belong to $D_{\rm perf}^{[0,3]}({}_{\OO_{X}}{\rm Mod})$. 
\end{prop}
\begin{proof} As in \cite[Proposition 9.7.2 (ii)]{nek-selmer}. 
\end{proof}

From the triangle  \eqref{ex tri} we extract an exact sequence
\beq\label{tilde ses}
\source{universal-p-adic-gross-zagier-formula}
0 \to H^{0}(G_{E, Sp}, M)\to \bigoplus_{w\in Sp} H^{0}(E_{w}, M^{-}_{w}) \to \wtil{H}^{1}_{f}(E, M)\to H^{1}_{f} (E, M)\to 0
\eeq
where the last term is the (Greenberg) Selmer group 
\beq\label{sel gr} H^{1}_{f} (E, M):= \Ker \left( H^{1}_{}(G_{E, Sp},M)\to \bigoplus_{w\in Sp}  H^{1} (E_{w}, M_{w}^{-})\right). 
\source{universal-p-adic-gross-zagier-formula}
\eeq

\subsubsection{Height pairings} \label{sec: ht}
\source{universal-p-adic-gross-zagier-formula}
For $?=\emptyset,\iota$, let  $M^{?}=(M^{?}, (M_{w}^{?,+}))$ be a strict Greenberg datum for $G_{E, Sp} $ over $X$. Suppose given a perfect  pairing of $\OO_{X}[G_{E, Sp}]$-modules
\beq \lb{skewp}
M\ot_{\OO_{X}}M^{\iota}\to \OO_{X}(1) \eeq
such that $M^{+}_{w}$ and $M^{+, \iota}_{w}$ are exact orthogonal of each other. Let $\Gamma_{F}$ be a profinite abelian group.


For every pair of  Greenberg data $M$, $M^{\iota}$ as above, there  is a   height pairing 
\beq \lb{hhh}
h_{M}\colon \wtil{H}^{1}_{f}(E, M)\ot_{\OO_{X}}  \wtil{H}^{1}_{f}(E, M^{\iota})
\to \OO_{X}\hat{\ot}\Gamma_{F}\eeq
constructed  in  \cite[\S11.1]{nek-selmer}. The following is a special case of \cite[Appendix C, Lemma 0.16]{venerucci-thesis}. 
\begin{prop}
\source{universal-p-adic-gross-zagier-formula}
\notready \lb{prop ht} 
For each regular point $x\in X$ and $P_{1}\ot P_{2}\in \wtil{H}^{1}_{f}(E,  M)\ot_{\OO_{X}}  \wtil{H}^{1}_{f}(E, M^{\iota})$, we have 
$$h_{M\ot \kappa(x)} (P_{1,x}, P_{2, x})=(h_{M}(P_{1},  P_{2}))(x).$$
\end{prop}


 Venerucci has defined height pairings in an even more general context. Let $M_{X}$ be a strict Greenberg datum over $X$ as above, let $Y\subset X$ be a  local complete intersection, and let $M_{Y}^{?}$ be the restriction of $M_{X}^{?}$.  Let $\mathscr{N}_{Y/X}^{*}$ be the conormal sheaf of $Y\to X$. 
Then there is a height pairing
\beq \lb{hnv} h_{M_{Y}/M_{X}}\colon  \wtil{H}^{1}_{f}(E, M_{Y})\ot_{\OO_{Y}}  \wtil{H}^{1}_{f}(E, M_{Y}^{\iota})\to \mathscr{N}_{Y/X}^{*}, \eeq
constructed in \cite[Appendix C, \S~0.21]{venerucci-thesis}. 

We  note its relation to \eqref{hhh} in a special case, and its symmetry properties in a conjugate-self-dual case. 


\begin{prop}
\source{universal-p-adic-gross-zagier-formula}
\notready \lb{prop ht2} The pairing \eqref{hnv}
satisfies the following properties. 
\begin{enumerate}
\item
Let $\Gamma_{F}$ be a profinite abelian quotient of $G_{E, Sp}$, let $X=Y\hat{\ts}_{\Spec \Q_{p}}\Spec \Z_{p}\llb \Gamma_{F}\rrb_{\Q_{p}}$ (where $\hat{\ts}= \eqref{hatprod}$), and assume that $M_{X}^{?}=M_{Y}^{?}\ot_{\Z_{p}} \Z_{p}\llb \Gamma_{F} \rrb$ for $?=\emptyset, \iota$, where if $?=\emptyset $  (respectively $?=\iota$) then $G_{E, Sp}$ acts on $\Gamma_{F}$ through the tautological character (respectively its inverse). 
Then $$h_{M_{Y}/M_{X}} = h_{M_{Y}}=\eqref{hhh}.$$
\item Suppose that there is an involution $\iota\colon X\to X$ stabilising $Y$ and such that $M_{X}^{\iota}= M_{X}\ot_{\OO_{X}, \iota}\OO_{X}$, $M_{X,w}^{+,\iota}= M_{X, w}^{+}\ot_{\OO_{X}, \iota}\OO_{X}$. 

 Let $\epsilon, \epsilon'\in\{\pm 1\}$. Assume that the pairing \eqref{skewp} is $\epsilon$-hermitian (\S~\ref{skew def}),
that  ${\rm d}_{Y/X}\iota= \epsilon'\id$  on $\mathscr{N}_{Y/X}^{*}$, and that there is an  $\OO_{Y}$-linear isomorphism 
$${\rm c}\colon \wtil{H}^{1}_{f}(E, M_{Y}^{\iota}) =  \wtil{H}^{1}_{f}(E, M_{Y})^{\iota} \to  \wtil{H}^{1}_{f}(E, M_{Y}) .$$
Then the pairing 
\beq\lb{skewsym}
h_{M_{Y}/M_{X}}^{\Box} \colon  \wtil{H}^{1}_{f}(E, M_{Y})\ot_{\OO_{Y}}  \wtil{H}^{1}_{f}(E, M_{Y}) &\to \mathscr{N}_{Y/X}^{*}\\
(z, z') &\mapsto h_{M_{Y}/M_{X}}(z, {\rm c} z')
 \eeq
 is $\epsilon\epsilon'$-symmetric. 
\end{enumerate}
\end{prop}
In our main application  in Theorem \ref{ugz thm}, we have  $Y=\X$ (or an open subset), the Hida family for $(\G\ts\H)'$; $X=\X^{\sharp}$, the Hida family for $\G\ts \H$ containing $X$; and  $M_{Y}=\cV$, $M_{X}=\cV^{\sharp}$, the corresponding universal $G_{E}$-representations. In that case, the height pairing $h_{\cV/\cV^{\sharp}}$ is simply $1/2 $ of the pairing $h_{\cV}$ of Proposition \ref{prop ht}. 



\begin{proof} Part 1 follows from the construction. (We omit further details since, by the remark preceding the present proof, we do not actually need it in this paper). 
We prove the symmetry properties. Let $I$ be the ideal sheaf of $Y$; up to restricting to some open subset of $X$ we may assume that $I$ is generated by a regular sequence $x=(x_{1}, \ldots, x_{r})$. Then $\calN_{Y/X}^* = I/I^2$ is finite locally free generated by $([x_1], \ldots, [x_r])$. Let $\partial_i \in N = (I/I^2)^\vee$ be the map $\partial_i([x_j])=\delta_{ij}$. It suffices to show that $h_i:=\partial_i \circ h$ is   $\epsilon \epsilon'$-hermitian 
 for all $i$. By \cite[Appendix C, Proposition 0.5]{venerucci-thesis}, $h_i$ is identified with $ h_{M_{Y}/M_{X_{i}}}$,  where $X_i=V_{X}((x_j)_{j\neq I})$ so that $Y=V_{X_{i}}(x_{i})$. Hence it suffices to prove the claim for $r=1$.

We argue similarly to \cite[Proof of Corollary 10.10]{venerucci-thesis}.  Assume thus $r=1$, write $x $ in place of $x_{1}$, and let $K$ be the fraction field of $X$. By  \cite{nek-selmer} and \cite[Appendix A, \S~ 0.7]{venerucci-thesis} we have an $\epsilon$-hermitian Cassels--Tate pairing 
$$\cup \colon \wtil{H}^{2}_{f}(E, M_{X})_{\text{$\OO_{X}$-tors}}
 \ot_{\OO_{X}} \wtil{H}^{2}_{f}(E, M_{X}^{\iota})_{\text{$\OO_{X}$-tors}} \to K/\OO_{X},$$
and by \cite[Appendix C, Proposition 0.17]{venerucci-thesis}, we have a map
 $$ i_{x}\colon \wtil{H}^{1}_{f}(E, M_{Y})\to  \wtil{H}^{2}_{f}(E, M_{X})[x]$$
 such that $h_{M_{Y}/M_{X}}$ coincides with 
$$ \wtil{H}^{1}_{f}(E, M_{Y}) \ot  \wtil{H}^{1}_{f}(E, M^{\iota}_{Y}) 
\stackrel{i_{x}\ot i_{x}^{\iota}}{\longrightarrow} 
 \wtil{H}^{2}_{f}(E, M_{X})[x]\ot  \wtil{H}^{2}_{f}(E, M^{\iota}_{X})[x]  
 \stackrel{\cup}{\longrightarrow}  x^{-1}\OO_{X}/\OO_{X}  \stackrel{[\cd x^{2}]}{\longrightarrow}  I/I^{2}= \calN^{*}_{Y/X}.$$
Since all the above maps are $\iota$-equivariant, we find that $h_{M_{Y}/M_{X}}$ is $\epsilon$-hermitian as well. The desired assertion follows from this and the fact that $\iota $ acts by $\eps'$ on $\calN^{*}_{Y/X}$.
 \end{proof}
